\section{Robustness of the Timing Results}\label{sec:robustness}

The answers to RQ1 and RQ2 depend on choices that a reader could question. We repeated the headline once for each such question, changing one thing at a time (Table~\ref{tab:robustness}).

\paragraph{Commits from one repository are not independent.} A repository that migrates several services at once adds several similar commits. Resampling whole repositories \citep{efron1993bootstrap}, \BootReps{} times over \BootRepos{} repositories, gives a 95\% interval of \AfterLateRepoLow--\AfterLateRepoHigh\% for the share after the shutdown and \ExactDaysRepoLow{} to \ExactDaysRepoHigh{} days for the median, which is almost the same as the per-commit intervals. Counting only the first migration of each repository away from each model gives \AfterFirst\%.

\paragraph{The label differs from Kim's.} Kim's rules (\texttt{migration}) give \AfterKimLabel\% on our data, the strict rules give \AfterStrict\%, and \texttt{migration\_away} without commits made with an AI coding tool gives \AfterNoAI\%. All of these are far from \KimAfterShutdown\%, so the difference from Kim does not come from the way we label migrations.

\paragraph{The label's mistakes may depend on timing.} If the label kept more false migrations on one side of the shutdown than the other, the share would be biased. On the hand-labelled sample, its precision is \PrecisionAfter\% for commits made after the shutdown (\PrecisionAfterLow--\PrecisionAfterHigh) and \PrecisionBefore\% for commits made before it (\PrecisionBeforeLow--\PrecisionBeforeHigh), so false migrations are more common before the shutdown. The migrations the label misses are also mostly early: \SampleMissedAfter\% of them happen after the shutdown. The two kinds of error therefore pull in opposite directions. In the sample, which covers all models, the label places \SampleLabelAfter\% of migrations after the shutdown; when each stratum is weighted by its size so that the missed migrations are counted too, the estimate for real migrations is \SampleTrueAfter\% (95\% CI \SampleTrueAfterLow--\SampleTrueAfterHigh, from resampling within strata). Because the sample is small, this check can only rule out a large bias.

\paragraph{The search may have missed commits.} On \CutDays{} days, GitHub returned only 1,000 of the results for an ID. \CutDaysBefore{} of these days fall before the model's shutdown, and together they contain only \CutRowsReweighted{} of the headline migrations. Counting each migration found on such a day as many times as the total reported by GitHub for that day divided by 1,000 leaves the share at \AfterReweighted\% (\AfterReweightedLow--\AfterReweightedHigh). The five IDs that were searched only from 30 days before their notice contribute \LateSearchedRows{} migrations to the headline, because none of their models had been shut down by \RetiredBy, except one of Google's.

\paragraph{OpenAI and Anthropic differ only by eye.} A two-sided Mann--Whitney U test \citep{mann1947}, with Cliff's $\delta$ \citep{cliff1993} as the effect size and named using the limits of \citet{romano2006}, finds no practical difference across all migrations ($\delta$~=~\CliffAll, \CliffAllSize, $p$~\CliffAllP) or across first migrations per repository ($\delta$~=~\CliffFirst, $p$~\CliffFirstP). Among migrations made after the shutdown, OpenAI projects migrate later ($\delta$~=~\CliffLate, \CliffLateSize, $p$~\CliffLateP).

\gentable{robustness}{tab:robustness}{Robustness of RQ1 and RQ2. Each line repeats the headline (OpenAI and Anthropic, models shut down by \RetiredBy) with one change. Intervals are 95\% Wilson intervals unless stated otherwise; the weighted lines use intervals from resampling repositories.}
